\section{Pairings, ordering, and comparison kernels}\label{S:nogo-pairings}

The second moment contains two ordered chains and all their Gaussian
pairings, including pairs within a chain. This structure offers several
natural reductions: estimate each pairing separately, compare with a single
chain, or copy the first-moment clock. The obstructions below distinguish
loss of scale from failure of an inequality. A majorant can be too large even
when every step used to obtain it is valid.

Throughout this section write $a=\alpha_0$ and
$\beta=4-\alpha-a$. Let $X_n$ denote the ordered Stratonovich term of degree
$n$, and put
\[
  a_k=\E X_{2k}(1,x),\qquad
  A_k=\sum_{m+n=2k}\E[X_m(1,x)X_n(1,x)],\qquad
  b_1=\E X_1(1,x)^2.
\]
Thus $A_1=2a_1+b_1$. Homogeneity supplies the factor $t^{\beta k}$.
The nonnegative convergent moment expansions used here are those of
Proposition~\ref{P:mixed}; they include unequal chain lengths. In dimension
one we use $G_s(z)=\frac12\mathbf1_{\{|z|<s\}}$.

\subsection{NG-01: one common estimate loses a factorial power}
\label{S:nogo-01}

Theorem~3.1 of~\cite{chen:author2026:time-dependency} controls all the
pairing patterns. Applying its homogeneous norm bound to each pattern is
therefore an attractive first attempt at the upper exponent. The issue is
what happens after optimization and summation.

\setcounter{nogo}{0}
\begin{nogo}[The common-parameter majorant has the wrong scale]
  \label{P:nogo-01}
  Under Hypothesis~\ref{H:homog}, let $P_k=(2k)!/(2^k k!)$ and $K>0$.
  The numerical majorant
  \[
    B_K(t)=1+\sum_{k\ge1}(2k+1)P_kK^k
    \inf_{\theta>0}e^{2\theta t}\theta^{-\beta k}
  \]
  satisfies $\log B_K(t)\asymp t^{\beta/(\beta-1)}$ as $t\to\infty$.
  Since $\beta/(\beta-1)>\kappa=\beta/(3-\alpha)$, this majorant cannot
  establish a finite upper constant on the scale $t^\kappa$.
\end{nogo}
\begin{proof}
  The infimum equals $e^{\beta k}(2t/(\beta k))^{\beta k}$.
  The inequalities $k!\le P_k\le2^k k!$ and
  $e^{-k}\le k!/k^k\le1$ show that $B_K$ lies between two series
  \[
    \sum_{k\ge0}\frac{(c t^\beta)^k}{(k!)^q},\qquad q=\beta-1>1,
  \]
  with different positive constants $c$. For $z>0$, this series is at most
  $\exp(qz^{1/q})$, by the $\ell^q$--$\ell^1$ inequality. Its term with
  $k=\lfloor z^{1/q}/2\rfloor$ is at least $2^{qk}$ for large $z$.
  This proves the two-sided growth. Finally, $\beta-1=3-\alpha-a<3-\alpha$.
  This proves No-go~\ref{P:nogo-01}.
\end{proof}

This restates the obstruction in Proposition~\ref{P:envelope} and
Remark~\ref{R:G0-verdict}. The lower estimate is for $B_K$, not for the
moment. The valid upper estimate improves the crude quadratic exponent,
but its factorial denominator misses $(k!)^a$. The split count $2k+1$
costs only an exponential factor in $k$; removing it does not repair the
loss. The same calculation for one chain, with $e^{\theta t}$, has the
same defect. Reopening requires an estimate retaining ordering or grouping
pairings before taking this common bound. Section~\ref{S:interval-upper}
provides such a different argument.

\subsection{NG-02: a single adjacent diagram is too large}
\label{S:nogo-02}

There are roughly an exponential factor times $k!$ pairings of $2k$
vertices. In white space, a uniform denominator $(k!)^3$ for each diagram
would therefore give the desired denominator $(k!)^2$ after summation.
The following example explains why that convenient sufficient estimate is
stronger than the problem permits.

\setcounter{nogo}{1}
\begin{nogo}[No uniform cubic factorial denominator per diagram]
  \label{P:nogo-02}
  Let $d=1$, $\Gamma=\delta_0$ and $0<a<1$. For the split $(0,2k)$
  and adjacent matching $\{(1,2),(3,4),\ldots,(2k-1,2k)\}$, its contribution is
  \[
    J_{\rm adj}(t)=
    \left[\frac{\Gamma_{\mathrm E}(1-a)}2\right]^k
    \frac{t^{(3-a)k}}{\Gamma_{\mathrm E}(1+(3-a)k)}.
  \]
  Consequently, there is no finite $C=C(a)$ bounding every diagram by
  $C^k t^{(3-a)k}/(k!)^3$ for all $k$ and $t>0$.
\end{nogo}
\begin{proof}
  Use root-forward increments $h_1,\ldots,h_{2k}>0$ with
  $\sum h_i<t$. Each adjacent contraction evaluates the even Green factor
  at zero, giving $1/2$, and integration of the preceding free increment
  gives mass $h_{2j-1}$. The temporal factor is $h_{2j}^{-a}$. Therefore,
  \[
    J_{\rm adj}(t)=2^{-k}
    \int_{\sum h_i<t}\prod_{j=1}^k h_{2j-1}h_{2j}^{-a}\,\ud h.
  \]
  Here delta contractions mean limits of Gaussian approximate identities.
  For completeness, fix $k$ and a strictly positive increment vector.
  After convolution the paired positions approach equality. The even Green
  factors are continuous at zero; the remaining indicator boundaries are
  finitely many hyperplanes in the $k$ free positions and have measure zero.
  The limit is thus the stated diagonal integral. All positions in the
  original chain lie in $[-t,t]$. Uniformly in the smoothing widths, replacing
  its Green indicators by one bounds the spatial integral by
  $2^{-2k}(2t)^k$: each disjoint smoothed pair has integral at most $2t$
  over $[-t,t]^2$. The temporal factors on the full time cube are a product
  of $k$ disjoint factors $|s-r|^{-a}$, each integrable because $a<1$.
  Dominated convergence justifies the fixed-order spatial limit through the
  time integration, and also temporal truncation limits.

  The Dirichlet simplex formula, with parameters $2$ and $1-a$ for each
  pair, now gives the asserted expression. Stirling's formula yields
  \[
    \log\big((k!)^3J_{\rm adj}(1)\big)=a k\log k+O(k).
  \]
  This exceeds $k\log C$ eventually for every fixed $C$.
  This proves No-go~\ref{P:nogo-02}.
\end{proof}

The adjacent matching is only one diagram. Its denominator is, up to
exponential factors, $(k!)^{3-a}$, which is stronger than the desired
summed denominator $(k!)^2$. Hence, this witness does not obstruct the
summed bound. A viable successor must retain the number of large diagrams,
or sum before estimating; it cannot assign the same cubic denominator to
all of them.

\subsection[NG-03: a squared-mean comparison fails]{NG-03: doubled variance does not give the proposed squared mean}
\label{S:nogo-03}

A positive mixture of exponentials of linear Gaussian functionals would
satisfy a useful comparison between the second moment and the mean at
noise amplitude $\sqrt2$. Such a representation must be established for
the wave rather than inferred from a parabolic analogy.

\setcounter{nogo}{2}
\begin{nogo}[Failure of the all-time comparison with unit prefactor]
  \label{P:nogo-03}
  For the admissible Riesz case $d=1$, $\Gamma(x)=|x|^{-1/2}$ and $a=1/2$,
  \[
    \E u(t,x)^2>\big(\E u_{\sqrt2}(t,x)\big)^2
  \]
  for all sufficiently small $t>0$, where $u_\lambda$ uses noise amplitude
  $\lambda$ and the same initial data.
\end{nogo}
\begin{proof}
  More generally, for $\Gamma(x)=c|x|^{-\alpha}$, $0<\alpha<1$,
  direct integration of the length-two chain gives
  \[
    a_1=\frac{c}{(1-\alpha)(2-\alpha-a)(3-\alpha-a)(4-\alpha-a)}.
  \]
  Indeed the spatial integral is
  $c h_1h_2^{1-\alpha}/(1-\alpha)$ and its temporal factor is $h_2^{-a}$;
  integrate over $h_1+h_2<1$. At $\alpha=a=1/2$, $c=1$, this gives $a_1=1/3$.
  For independent uniforms $U$ on $[-s,s]$ and $V$ on $[-r,r]$, $0<r\le s$,
  \[
    \E|U-V|=\frac{s}{2}+\frac{r^2}{6s}\le\frac{2s}{3}.
  \]
  Jensen's inequality for $z\mapsto z^{-\alpha}$ implies
  \[
    \int G_s(x)G_r(y)|x-y|^{-\alpha}\,\ud x\ud y
    \ge (3/2)^\alpha s^{1-\alpha}r.
  \]
  Integrating on the two temporal triangles gives
  \[
    b_1\ge
    \frac{2(3/2)^\alpha}{(1-a)(2-a)(4-\alpha-a)}
    =\frac89\sqrt{\frac32}>\frac23=2a_1.
  \]
  The comparison is a finite-time failure, not just a formal coefficient test.
  Put $z=t^\beta$. The two nonnegative series under comparison have expansions
  \[
    \E u(t,x)^2=1+(2a_1+b_1)z+O(z^2),\qquad
    (\E u_{\sqrt2}(t,x))^2=1+4a_1z+O(z^2).
  \]
  To justify the remainders, choose any fixed positive $z_0$ at which these
  series converge, using the moment expansion also at amplitude $\sqrt2$.
  For each nonnegative coefficient sequence $(c_j)$,
  $\sum_{j\ge2}c_jz^j\le(z/z_0)^2\sum_{j\ge2}c_jz_0^j$ when
  $0\le z\le z_0$. Apply this to both series, including the nonnegative
  Cauchy product for the squared mean. Their difference therefore has the
  strictly positive leading coefficient $b_1-2a_1$.
  This proves No-go~\ref{P:nogo-03}.
\end{proof}

For a deterministic nonnegative mixture of Gaussian exponentials, the
inequality follows by Tonelli and
$\|h+k\|^2/2\le\|h\|^2+\|k\|^2$, provided the same mixture represents
all amplitudes in the comparison. That proposed mechanism is excluded by
the example. White noise, an eventual large-time comparison, or a prefactor
larger than one is not excluded. Reopening requires a restriction or a new
comparison that removes the contradicted first-coefficient requirement;
it cannot merely assert this positive-mixture representation.

\subsection[NG-04: constants near the time boundary]{NG-04: a coefficient constant cannot be uniform at the time boundary}
\label{S:nogo-04}

Comparing $A_k$ with $a_k$ retains the self-pairings of the Stratonovich
problem and would transfer a first-moment bound to the second moment.
However, a prescribed asymptotic constant must not be imposed on every
low-order coefficient without checking parameter uniformity.

\setcounter{nogo}{3}
\begin{nogo}[Divergence of the first coefficient ratio]
  \label{P:nogo-04}
  Fix $0<\alpha<1$ and $c>0$, with $\Gamma(x)=c|x|^{-\alpha}$ in $d=1$.
  Then
  \[
    \lim_{a\uparrow1}(1-a)\frac{A_1}{a_1}
    =2^{2-\alpha}(1-\alpha)>0.
  \]
  In particular no constant depending on $\alpha$ alone, including
  $2^{4-\alpha}$, can satisfy $A_k\le C^k a_k$ for every $k\ge1$
  and every $0<a<1$ at this fixed spatial covariance.
\end{nogo}
\begin{proof}
  Let
  $K(s,r)=\int G_s(x)G_r(y)c|x-y|^{-\alpha}\,\ud x\ud y$.
  This is continuous on $[0,1]^2$ by local integrability of the Riesz kernel
  and dominated convergence on a fixed square. Direct integration gives
  \[
    K(s,s)=\frac{c\,2^{1-\alpha}s^{2-\alpha}}
    {(1-\alpha)(2-\alpha)}.
  \]
  Writing $v=s-r$, define the continuous function
  $H(v)=2\int_v^1 K(s,s-v)\,\ud s$. Then
  \[
    b_1=\int_0^1v^{-a}H(v)\,\ud v,\qquad
    H(0)=\frac{c\,2^{2-\alpha}}
    {(1-\alpha)(2-\alpha)(3-\alpha)}.
  \]
  For any continuous $H$ on $[0,1]$,
  $\varepsilon\int_0^1v^{-1+\varepsilon}H(v)\,\ud v\to H(0)$.
  Indeed the measures have mass one, and their mass on $[\delta,1]$ is
  $1-\delta^\varepsilon\to0$; continuity controls $[0,\delta]$.
  Thus $(1-a)b_1\to H(0)$. The formula for $a_1$ in the preceding proof gives
  \[
    a_1\longrightarrow
    \frac{c}{(1-\alpha)^2(2-\alpha)(3-\alpha)}.
  \]
  Combine these limits with $A_1=2a_1+b_1$.
  This proves No-go~\ref{P:nogo-04}.
\end{proof}

For each fixed $a<1$, a finite constant $C=C(\alpha,a)$ remains possible;
the argument only forces it to diverge as $a\uparrow1$. It also says
nothing against comparisons beginning at large order or large time with
suitable prefactors. Reopening requires explicitly allowing that parameter
dependence or removing the low orders. Historical static-endpoint tests
likewise constrain an endpoint or a uniform family, not a separately
chosen constant at an arbitrary positive $a$.

\subsection{NG-05: chronological merging breaks wave support}
\label{S:nogo-05}

A single chain has the monotone ordering needed by the first-moment
representation. Sorting two chains by time looks like a direct way to
recover it. Finite propagation prevents the resulting pointwise comparison.

\setcounter{nogo}{4}
\begin{nogo}[Support obstruction to pointwise concatenation]
  \label{P:nogo-05}
  In $d=1$, no finite constant can bound every product of two root-forward
  chain kernels by the single-chain kernel obtained by chronological merging.
  The failure occurs in the $(2,2)$ split on a set of positive measure for
  Riesz covariance and, with domestic pairs, for white covariance.
\end{nogo}
\begin{proof}
  Work at horizon one with root $(0,0)$. The first chain has times
  $(1/4,3/4)$ and positions $(x,x)$, $x=1/5$; the second has times
  $(1/2,7/8)$ and positions $(y,y)$, $y=-2/5$.
  The four original displacements and elapsed times are
  \[
    (1/5,1/4),\quad(0,1/2),\quad(2/5,1/2),\quad(0,3/8).
  \]
  Each displacement is strictly smaller than its elapsed time, so their
  Green product is $2^{-4}$. The merged chain instead includes the jump
  from $(1/4,1/5)$ to $(1/2,-2/5)$. Its displacement is $3/5>1/4$,
  so that Green factor vanishes. All these inequalities are strict.

  For white noise choose domestic pairs within each original chain. The
  delta contractions leave precisely the two free spatial variables $x,y$.
  Open neighborhoods of the displayed four times and these two variables
  preserve both conclusions and have positive contracted measure. Gaussian
  regularization justifies this contraction as in No-go~\ref{P:nogo-02}.
  For Riesz covariance perturb all four spatial positions independently
  in small neighborhoods. Avoiding the domestic diagonals gives a set of
  positive ordinary measure on which the covariance factors are finite and
  strictly positive, while the same support defect persists.
  No multiplier can turn the zero merged product into a positive upper bound.
  This proves No-go~\ref{P:nogo-05}.
\end{proof}

\begin{figure}[htbp]
  \centering
  \begin{tikzpicture}[x=5cm,y=5cm,>=stealth,font=\small]
    \draw[->,gray] (-.65,0) -- (.55,0) node[right] {$x$};
    \draw[->,gray] (0,-.04) -- (0,1.04) node[above] {$s$};
    \draw[gray!35,densely dotted] (-.6,.6) -- (0,0) -- (.5,.5);
    \coordinate (O) at (0,0);
    \coordinate (A) at (.2,.25);
    \coordinate (B) at (.2,.75);
    \coordinate (C) at (-.4,.5);
    \coordinate (D) at (-.4,.875);
    \draw[hamblue,very thick,->] (O) -- (A);
    \draw[hamblue,very thick,->] (A) -- (B);
    \draw[hamteal,very thick,->] (O) -- (C);
    \draw[hamteal,very thick,->] (C) -- (D);
    \draw[hamgold,very thick,dashed,->] (A) -- (C)
    node[midway,above,sloped,fill=white,inner sep=2pt]
    {$3/5>1/4$};
    \fill (O) circle (.6pt) node[below right] {root $(0,0)$};
    \fill[hamblue] (A) circle (.7pt)
    node[right,black] {$(1/5,1/4)$};
    \fill[hamblue] (B) circle (.7pt)
    node[right,black] {$(1/5,3/4)$};
    \fill[hamteal] (C) circle (.7pt)
    node[left,black] {$(-2/5,1/2)$};
    \fill[hamteal] (D) circle (.7pt)
    node[left,black] {$(-2/5,7/8)$};
    \node[hamblue,anchor=west] at (.65,.94) {first chain};
    \node[hamteal,anchor=west] at (.65,.84) {second chain};
    \node[hamgold,anchor=west] at (.65,.74) {merged jump};
  \end{tikzpicture}
  \caption{Root-forward coordinates $(x,s)$ for No-go~\ref{P:nogo-05}.
    Each solid edge lies strictly inside its propagation cone. Chronological
    merging inserts the dashed jump, whose distance $3/5$ exceeds its elapsed
    time $1/4$. The picture refutes pointwise concatenation, not an integrated
    comparison. For domestic white pairs the two free positions are \mbox{$x$ and $y$}.}
  \label{F:nogo-chain-geometry}
\end{figure}


The time coordinates here are root-forward chain coordinates; applying a
terminal-to-point convention again would change the object being tested.
The result concerns pointwise concatenation. Integrated comparisons and
comparisons of grouped Wick sums remain possible. Reopening requires
averaging, extra displacement variables, or another mechanism that preserves
the two original propagation constraints instead of imposing the false
merged one.

\subsection{NG-06: erasing geometry before summation loses the scale}
\label{S:nogo-06}

After white contractions, the temporal pairing sum is symmetric. One can
exploit that symmetry immediately by forgetting the spatial indicators.
The resulting bound is easy to sum but too weak for the desired exponent.

\setcounter{nogo}{5}
\begin{nogo}[The geometry-free majorant]
  \label{P:nogo-06}
  For $d=1$, $\Gamma=\delta_0$, $0<a<1$, replacing every wave indicator by
  one while retaining only the spatial box $[-t,t]$ gives the coefficient
  majorant
  \[
    A_k t^{(3-a)k}\le
    \frac{1}{k!}\left[\frac{2t^{3-a}}{(1-a)(2-a)}\right]^k.
  \]
  Its sum is $\exp(2t^{3-a}/((1-a)(2-a)))$, whose time exponent
  $3-a$ exceeds the target $(3-a)/2$.
\end{nogo}
\begin{proof}
  The spatial replacement gives $(2t)^k/2^{2k}=t^k/2^k$ for each pairing.
  After this replacement the sum over all pairings is symmetric in the
  $2k$ times. Its integral over a split $(m,n)$ simplex equals the cube
  integral divided by $m!n!$. The identities
  \[
    \sum_{m+n=2k}\frac1{m!n!}=\frac{2^{2k}}{(2k)!},\qquad
    P_k=\frac{(2k)!}{2^k k!},\qquad
    \int_0^t\int_0^t|s-r|^{-a}\,\ud s\ud r
    =\frac{2t^{2-a}}{(1-a)(2-a)}
  \]
  give the stated product and its exponential sum.
  This proves No-go~\ref{P:nogo-06}.
\end{proof}

This is a valid upper bound and a failure of its scale, not a moment lower
bound. Symmetry itself remains useful. The reopening condition is a spatial
estimate with enough order dependence retained before temporal
symmetrization; Proposition~\ref{P:d1-upper} retains this order dependence
and sums the temporal pairings afterwards (see
Remark~\ref{R:d1-upper-checks}).

\subsection{NG-07: block labels are not chronological ranks}
\label{S:nogo-07}

Temporal singularities favor close paired times. A proposed penalty for
long pairs must therefore measure time separation or actual chronological
rank. Labeling all vertices of one chain before those of the other creates
an artificial distance unrelated to either.

\setcounter{nogo}{6}
\begin{nogo}[Large label span does not imply a cubic denominator]
  \label{P:nogo-07}
  In white space with $0<a<1$, consider two length-$k$ chains and the matching
  $j\leftrightarrow k+j$. Although every pair has label distance $k$, its
  unit-time diagram satisfies
  \[
    J_k(1)\ge \left(\frac{3\,4^a}{128}\right)^k k^{-(3-a)k}.
  \]
  Thus no estimate assigning $C^k/(k!)^3$ to these diagrams solely on account
  of their block-label span can hold uniformly in $k$.
\end{nogo}
\begin{proof}
  Let the chain times be $s_j,r_j$, with increments $h_j,g_j$ from the
  common root. White contractions identify their paired positions, and
  integration of the successive spatial increments gives
  $2^{-k}\prod_{j=1}^k\min(h_j,g_j)$. Restrict both $s_j,r_j$ to
  $((j-1/4)/k,j/k)$. The domain has volume $(4k)^{-2k}$,
  each increment is at least $3/(4k)$, and each paired difference is at most
  $1/(4k)$. Multiplication gives the lower bound. The contraction can again
  be obtained at fixed $k$ by Gaussian smoothing; the spatial box bound and
  integrable disjoint temporal factors justify the limit.
  Stirling's formula implies
  $\log((k!)^3J_k(1))\ge a k\log k+O(k)$, contradicting any fixed $C$.
  This proves No-go~\ref{P:nogo-07}.
\end{proof}

% [held back: collaborators' unpublished work]

\subsection[NG-10: crossing covariance survives merging]{NG-10: an unweighted merge does not remove a crossing covariance}
\label{S:nogo-10}

A two-wave convolution estimate can remove an interior chain vertex when
no other factor depends on it. In a crossing pairing that vertex still
interacts with a remote partner. The following order-four calculation
identifies the missing weight before any all-order iteration is attempted.

\setcounter{nogo}{9}
\begin{nogo}[No finite spatial-supremum pullout in the merge lens]
  \label{P:nogo-10}
  In $d=2$, take $\Gamma(z)=c|z|^{-\alpha}$, $c>0$,
  $0<\alpha_0<1/2$ and $3/2\le\alpha<2-\alpha_0$.
  For the root-forward $(0,4)$ diagram
  with pairs $(1,3),(2,4)$, merging vertex two requires the weighted integral
  \[
    W=c\int_a^b(T-v)^{-\alpha_0}
    \int_{\mathbb R^2}G_{v-a}(y-x)G_{b-v}(z-y)|w-y|^{-\alpha}
    \,\ud y\ud v,
    \qquad 0<a<b<T.
  \]
  Here the letters $a,b,T$ are vertex times, so the temporal exponent is
  written $\alpha_0$. The spatial weight has no finite uniform essential
  supremum on admissible merge lenses. Hence, pulling it out by a uniform
  supremum cannot reduce $W$ to an unweighted two-wave convolution.
\end{nogo}
\begin{proof}
  Set $t=5$, $a=1$, $b=3$, $T=4$, $x=(0,0)$,
  $z=(1/2,0)$ and $w=(1/4,0)$. For $v$ near $2$, the point $y=w$
  is strictly inside both disks $|y-x|<v-a$ and $|z-y|<b-v$.
  It also satisfies the outer constraint $|w-z|<T-b$, and $|x|<a$.
  There is therefore an interior ball about $w$ in the merge lens.
  For every $M>0$ a smaller ball has positive Lebesgue measure and
  $|w-y|^{-\alpha}>M$ there. This proves infinite essential supremum;
  the strict inequalities persist under small changes of the fixed data.
  The other pair contributes $(b-a)^{-\alpha_0}c|z-x|^{-\alpha}$ outside
  the merge, whereas the two factors for pair $(2,4)$ depend on $v,y$
  and must remain in $W$. The temporal weight alone admits the bound
  $(T-v)^{-\alpha_0}\le(T-b)^{-\alpha_0}$.
  This proves No-go~\ref{P:nogo-10}.
\end{proof}

The unweighted estimate of~\cite{bolanos-guerrero.nualart.ea:21:averaging}
cannot be substituted for this different object without controlling the
weight. Infinite essential supremum does not imply divergence of $W$ or
of the diagram: the Riesz singularity is locally integrable for $\alpha<2$.
A reopening condition is a covariance-weighted convolution inequality,
or an explicit envelope integrable against the remaining diagram factors,
with the remote-gap cost retained. No such inequality is asserted here.

\subsection{NG-17: direct copying of the first-moment clock undercounts}
\label{S:nogo-17}

The source's first-moment clock representation suggests attaching one
Brownian clock to each ordered chain. A degree-two cross pair already tests
whether this copying preserves the interaction. It does not, even when
the two clocks are shared.

\setcounter{nogo}{16}
\begin{nogo}[Constant-correlation clock copying misses a positive factor]
  \label{P:nogo-17}
  Under the nonzero covariance hypotheses of
  Proposition~\ref{P:clock-copy-obstruction}, let $K_\Gamma(\theta)$ be the
  true cross-pair transform and $J_{\Gamma,\rho}(\theta)$ the direct
  unweighted candidate~\eqref{E:naive-clock-candidate}, $-1\le\rho\le1$.
  Then for every $\theta>0$,
  \[
    0<J_{\Gamma,\rho}(\theta)\le J_{\Gamma,1}(\theta)<K_\Gamma(\theta).
  \]
  If $\alpha=1$, then $J_{\Gamma,1}=(2-2^a)K_\Gamma$.
\end{nogo}
\begin{proof}
  Recall $\mathcal R_\Gamma(q)=\int(q^2+|\xi|^2)^{-1}\mu(\ud\xi)$.
  The exact transform identities proved in
  Proposition~\ref{P:clock-copy-obstruction} give
  \[
    w_\theta(\lambda)=
    \frac{\lambda^{a-2}}{2\theta\Gamma_{\mathrm E}(a)}
    [\mathcal R_\Gamma(\theta)-\mathcal R_\Gamma(\theta+\lambda)],
    \qquad
    K_\Gamma=\int_0^\infty w_\theta(\lambda)\,\ud\lambda,
  \]
  \[
    J_{\Gamma,1}=\int_0^\infty
    \frac{2\theta}{2\theta+\lambda}w_\theta(\lambda)\,\ud\lambda.
  \]
  The density is strictly positive because $\mu\ne0$ and the resolvent
  strictly decreases. It is integrable: homogeneity bounds it by multiples
  of $\lambda^{a-1}$ near zero and $\lambda^{a-2}$ at infinity.
  The multiplier is strictly between zero and one. For smaller correlations,
  $u+v-2\rho\min(u,v)\ge|u-v|$, and the positive Laplace integral
  $\Psi_a(r,v)$ in~\eqref{E:averaged-clock-kernel} decreases with $v$.
  This proves the chain of inequalities. When $\alpha=1$, insert
  $\mathcal R_\Gamma(q)=\mathcal R_\Gamma(1)/q$ and use
  $\int_0^\infty x^{a-1}/(x+c)\,\ud x
  =c^{a-1}\Gamma_{\mathrm E}(a)\Gamma_{\mathrm E}(1-a)$ to obtain the ratio.
  This proves No-go~\ref{P:nogo-17}.
\end{proof}

This is the already proved book obstruction, not a new clock impossibility
theorem. At the static endpoint the discrepancy disappears, so homogeneity
and static checks alone cannot detect it. At degree two the missing
shared-clock multiplier is exactly $1+\lambda/(2\theta)$; no higher-order
correction rule follows. Reopening requires relative-order weights or a
different coupling or representation that reproduces the actual cross-pair
transform. Remark~\ref{R:clock-copy-scope} gives the full boundary of this
statement, including the distinction between gap transforms and physical
unequal-time moments.

\subsection[NG-18: an oversized middle-cross envelope]{NG-18: erasing the middle-cross shift makes an envelope too large}
\label{S:nogo-18}

A later white-space comparison retains four vertices and a middle cross
pair. Removing the clock-overlap penalty gives a positive upper bound.
One further simplification, erasing the shift of its root kernel, also
increases the bound, but increases it beyond the comparator at an explicit
parameter point. This is an obstruction to closing that envelope, not a
counterexample to the original comparison.

\setcounter{nogo}{17}
\begin{nogo}[The shift-erased envelope is not below the wave transform]
  \label{P:nogo-18}
  Consider the white $(1,3)$ diagram with middle-cross partition $02\cdot13$,
  common gap damping one, cross rate $\lambda\ge0$ and self rate $\nu\ge0$.
  Put $b=1+\nu$, $c=1+\lambda$, $A=\lambda+\lambda^2/2$. Its actual wave
  transform $I$ and the shift-erased clock upper envelope $M_b(0;A)$ satisfy
  at $b=1,c=11/10$
  \[
    I=\frac{545}{11616},\qquad M_b(0;A)=\frac{1025}{21483},\qquad
    I-M_b(0;A)=-\frac{6005}{7562016}<0.
  \]
  Consequently the proposed pointwise closure $M_b(0;A)\le I$ is false.
\end{nogo}
\begin{proof}
  We give the integral identifications so that the rational test specifies
  which bound fails. With singleton gap $u$ and long-chain gaps $v,w,z$,
  the spatial factor is
  $\int G_u(X)G_v(Y)G_w(X-Y)G_z(Y-X)\,\ud X\ud Y$.
  Set $R=u+v$, $q=u-v$. Integration in $R$, including its Jacobian,
  gives the exact wave expression
  \[
    I=\frac1{16}\int_{w,z>0}\int_{\mathbb R}
    e^{-b(w+z)-\lambda|q-w|}
    \int_0^{\min(w,z)}e^{-\max(|q|,t)}\,\ud t\ud q\ud w\ud z.
  \]
  To check the spatial reduction, the derivative of the area
  $\{|X|<u,|Y|<v,|X-Y|<C\}$ with respect to $C$ is
  $2(R-\max(C,|q|))_+$; the area vanishes at $C=0$.
  Integrate this against $e^{-R}/32$ for $R>|q|$.
  Next integrate $z\ge t$ and put $w=t+x$. Splitting the remaining $q$
  integral into $q<0$, $0\le q\le t$, and $q\ge t$ gives respectively
  \[
    \frac{1/c+1/(2b+\lambda+1)}{16b(b+\lambda)(2b+2\lambda+1)},\quad
    \frac1{16b(b+\lambda)(2b+1)(2b+1+\lambda)},
  \]
  \[
    \frac{1/(b+\lambda)+1/c}{16b(2b+1)(b+1)}.
  \]
  Adding gives
  \[
    I=\frac{2b^2+4bc+2c^2-c+1}
    {8bc(b+1)(2b+1)(b+c-1)(2b+2c-1)}.
  \]

  For the clock envelope write $h(r)=(2\pi r)^{-1/2}$ and
  $m(s)=1/(\pi\sqrt2\sqrt{s-1/2})$, $s>1/2$.
  The two spatial pair differences have covariance matrix
  $\left(\begin{smallmatrix}u+v+w&-w\\-w&w+z\end{smallmatrix}\right)$.
  Writing $r=w+z$, its determinant is $rR+wz$, so their joint white
  density factors as $h(r)h(R+wz/r)$. After deleting the nonnegative
  clock-overlap penalty, the root kernel is
  \[
    H_\gamma(q)=\int_{|q|}^\infty e^{-R/2}h(R+\gamma)\,\ud R
    =\int_{1/2}^\infty\frac{m(s)}s
    e^{-(s-1/2)\gamma}e^{-s|q|}\,\ud s,
    \qquad \gamma=\frac{wz}{r}.
  \]
  The equality follows from the positive Gamma integral for $h$.
  Its convolution
  with $e^{-A|q-w|}$ equals $2S_{A,s}(w)/(s+A)$, where
  \[
    S_{A,s}(w)=e^{-Aw}+\int_0^w A e^{-At}e^{-s(w-t)}\,\ud t\le1.
  \]
  Replacing $S$ by one is precisely the shift erasure. More generally,
  replacing it by $e^{-\eta w}$ defines
  \[
    M_b(\eta;A)=\frac1{16}\int_{1/2}^\infty
    \frac{m(s)}{s(s+A)}\int_0^\infty h(r)e^{-b^2r/2}
    \int_0^r e^{-(s-1/2)w(r-w)/r-\eta w}\,\ud w\ud r\ud s.
  \]
  All integrals here are positive. Put $w=rt$ and $p=\sqrt{b^2+2\eta}$.
  The elementary bridge integral is
  \[
    \int_0^1[b^2+2\eta t+2\tau t(1-t)]^{-3/2}\,\ud t
    =\frac{2(p+b)}{bp((p+b)^2+2\tau)}.
  \]
  It follows by integrating a quadratic power; at $\eta=\tau=0$ both sides
  are $b^{-3}$. Substitution $z=\sqrt{2s-1}$ and partial fractions give
  \begin{align*}
    M_b(\eta;A)&=\frac{p+b}{2\pi bp}
    \int_0^\infty\frac{\ud z}
    {(z^2+1)(z^2+c^2)(z^2+(p+b)^2)}\\
    &=\frac{1+c+p+b}{4bp c(c+1)(c+p+b)(1+p+b)}.
  \end{align*}
  The usual identity $\int_0^\infty(z^2+q^2)^{-1}\ud z=\pi/(2q)$
  proves the partial fraction evaluation; coincident rates follow by
  continuity. At $\eta=0$, $p=b$. Substituting the asserted rational
  parameters in these two expressions gives the displayed strict gap.
  This proves No-go~\ref{P:nogo-18}.
\end{proof}

The boundary value $\nu=0$ is a legitimate fixed-parameter transform;
continuity of both rational expressions also gives an overshoot for some
$\nu>0$ nearby. The no-go does not refute $J\le I$, where $J$ is the
original shared-clock candidate, nor does it decide an integrated
comparison using this envelope. The surviving middle-cross comparison
retains $S_{A,1/2}(w)$, instead of replacing it by one, and obtains a smaller
envelope. Reopening this particular closure requires keeping a shift cost
or proving a different integrated estimate. The historical standalone
record was not separately adjudicated; the subsequent accepted
middle-cross analysis and this reconstructed proof do not change that
historical review status.
